\documentclass[11pt,a4paper]{article}

\usepackage[margin=1in]{geometry}
\usepackage{amsmath}
\usepackage{amssymb}
\usepackage{booktabs}
\usepackage{array}
\usepackage{hyperref}

\hypersetup{
  colorlinks=true,
  linkcolor=blue,
  urlcolor=blue
}

\title{LCMS-to-LCMS Transformation for Resonance Contributions\\to the 3D Femtoscopic Source}
\author{}
\date{April 29, 2026}

\begin{document}

\maketitle

\section{Scope}

The goal is to model the transformation of the core separation vector from the
source-frame LCMS to the final daughter-pair LCMS. The experimental 3D
femtoscopic source to be fitted is defined in the daughter-pair LCMS, with
coordinates
\[
  (r_{\mathrm{out}}, r_{\mathrm{side}}, r_{\mathrm{long}}).
\]
Therefore no transformation to the pair rest frame is required here.

The relevant transformation is
\[
  \mathrm{LCMS}_{\mathrm{core/source}}
  \rightarrow
  \mathrm{LCMS}_{\mathrm{daughter}},
\]
followed by the addition of resonance decay propagation shifts when a particle
is produced from a resonance.

\section{Coordinate Convention}

For a pair of particles with total four-momentum
\[
  P = p_1 + p_2,
\]
the LCMS is defined by removing the longitudinal pair velocity,
\[
  \beta_z = \frac{P_z}{P_0}.
\]
The outward axis is defined by the transverse pair momentum,
\[
  \hat{o} =
  \frac{(P_x, P_y, 0)}
       {\sqrt{P_x^2 + P_y^2}},
\]
the sideward axis by
\[
  \hat{s} = \hat{z} \times \hat{o},
\]
and the longitudinal axis by
\[
  \hat{l} = \hat{z}.
\]

\section{Core Separation in the Source LCMS}

The sampled core separation is assumed to be generated in the source-pair LCMS:
\[
  \Delta x_S =
  \left(
    \Delta t_S,
    \Delta \vec{r}_S
  \right),
  \qquad
  \Delta t_S = 0,
\]
with
\[
  \Delta \vec{r}_S =
  r_{\mathrm{out}}^S \hat{o}_S
  + r_{\mathrm{side}}^S \hat{s}_S
  + r_{\mathrm{long}}^S \hat{z}.
\]
The superscript \(S\) denotes the source-pair LCMS, not necessarily a
resonance-resonance frame. Its exact definition depends on the topology.

\section{Source LCMS to Daughter LCMS}

Let the source-pair LCMS have longitudinal velocity
\[
  \beta_S = \frac{P_{z,S}}{P_{0,S}},
  \qquad
  \gamma_S = \frac{1}{\sqrt{1-\beta_S^2}},
\]
and the daughter-pair LCMS have
\[
  \beta_D = \frac{P_{z,D}}{P_{0,D}},
  \qquad
  \gamma_D = \frac{1}{\sqrt{1-\beta_D^2}}.
\]

Starting from \(\Delta t_S = 0\), the transverse part of the separation is
unchanged by the longitudinal boost, while the longitudinal coordinate and time
are mixed. With the convention used in the EPOS helper macro,
\[
  p_z^{\mathrm{LCMS}} =
  \gamma (p_z - \beta E),
\]
the source-LCMS separation transforms into the daughter LCMS as
\[
  \Delta z_D =
  \gamma_S \gamma_D
  (1 - \beta_S \beta_D)
  \Delta z_S,
\]
and
\[
  \Delta t_D =
  \gamma_S \gamma_D
  (\beta_S - \beta_D)
  \Delta z_S.
\]
Equivalently, in terms of pair rapidities \(Y_S\) and \(Y_D\),
\[
  \Delta z_D =
  \cosh(Y_S - Y_D)\Delta z_S,
  \qquad
  \Delta t_D =
  \sinh(Y_S - Y_D)\Delta z_S.
\]

Since the target source is defined in the daughter LCMS, the relevant spatial
coordinates are obtained by projecting the transverse vector onto the daughter
OSL axes:
\[
  r_{\mathrm{out}}^D =
  \Delta \vec{r}_{T,S} \cdot \hat{o}_D,
\]
\[
  r_{\mathrm{side}}^D =
  \Delta \vec{r}_{T,S} \cdot \hat{s}_D,
\]
\[
  r_{\mathrm{long}}^D =
  \Delta z_D.
\]
Thus the out-side part is mainly a transverse OSL-axis rotation, while the
longitudinal part is controlled by the rapidity difference between the two LCMS
frames.

\section{Topology Dependence}

The boost formula is universal. What changes between resonance topologies is
the definition of the source pair, and therefore the source LCMS itself.

\begin{center}
\begin{tabular}{>{\raggedright\arraybackslash}p{0.12\linewidth}
                >{\raggedright\arraybackslash}p{0.34\linewidth}
                >{\raggedright\arraybackslash}p{0.34\linewidth}}
\toprule
Topology & Source pair defining \(\mathrm{LCMS}_S\) & Final pair defining \(\mathrm{LCMS}_D\) \\
\midrule
RR &
\((R_1, R_2)\) &
\((\pi_1, \pi_2)\) \\
PR &
\((\pi_1^{\mathrm{prim}}, R_2)\) &
\((\pi_1^{\mathrm{prim}}, \pi_2)\) \\
RP &
\((R_1, \pi_2^{\mathrm{prim}})\) &
\((\pi_1, \pi_2^{\mathrm{prim}})\) \\
\bottomrule
\end{tabular}
\end{center}

Therefore the EPOS information should distinguish the three cases at least by a
topology flag. The transformation can still be implemented by one common
function:
\[
  \mathrm{LCMS}(p_{1,\mathrm{source}} + p_{2,\mathrm{source}})
  \rightarrow
  \mathrm{LCMS}(p_{1,\mathrm{daughter}} + p_{2,\mathrm{daughter}}).
\]

\section{Recommended EPOS Storage}

For each sampled entry, the robust storage choice is to keep the four lab-frame
four-momenta needed to reconstruct both pair frames:
\[
  p_{1,\mathrm{source}},
  \quad
  p_{2,\mathrm{source}},
  \quad
  p_{1,\mathrm{daughter}},
  \quad
  p_{2,\mathrm{daughter}},
\]
together with a topology label:
\[
  \mathrm{topology} \in \{\mathrm{RR}, \mathrm{PR}, \mathrm{RP}\}.
\]

The mapping is
\[
\begin{array}{lll}
\mathrm{RR}: &
p_{1,\mathrm{source}} = R_1, &
p_{2,\mathrm{source}} = R_2, \\
\mathrm{PR}: &
p_{1,\mathrm{source}} = \pi_1^{\mathrm{prim}}, &
p_{2,\mathrm{source}} = R_2, \\
\mathrm{RP}: &
p_{1,\mathrm{source}} = R_1, &
p_{2,\mathrm{source}} = \pi_2^{\mathrm{prim}}.
\end{array}
\]
The daughter pair is always the final femtoscopic pair.

Storing only the mother momentum already boosted into the daughter LCMS is not
sufficient for this coordinate transformation, because the full transformation
requires reconstructing both the source-pair OSL axes and the daughter-pair OSL
axes, as well as the rapidity difference between the two LCMS frames.

\section{Practical Consequence for \texttt{DLM\_CleverReso3D}}

The implementation should avoid hard-coding separate boost formulas for RR, PR,
and RP. Instead, topology should determine which four-momenta are passed as the
source-pair momenta. The transformation routine can then be common to all
topologies:
\[
  (r_{\mathrm{out}}^S,r_{\mathrm{side}}^S,r_{\mathrm{long}}^S)
  \rightarrow
  (r_{\mathrm{out}}^D,r_{\mathrm{side}}^D,r_{\mathrm{long}}^D).
\]

The subsequent decay contribution should preserve the particle ordering. For a
relative coordinate convention \(r = x_1 - x_2\), the resonance propagation
shift enters with opposite signs for particle 1 and particle 2. This is why PR
and RP should not be merged without retaining the particle labels.

\end{document}
